% ============================================================
% Related Work  (extended version)
% Spine: language invariance / cross-lingual consistency.
% Uses plain \cite{} so it is compatible with any bibliography
% backend (natbib, biblatex, or the default BibTeX). All keys
% are defined in references.bib.
% ============================================================

\section{Related Work}
\label{sec:related-work}

This section reviews the work most relevant to the thesis: LLM-based ranking,
multilingual fact-checking, cross-lingual consistency, numerical reasoning, and
preference-based fine-tuning. The central question is whether language invariance
can be useful not only as something to measure, but also as something to optimize
for: if a model is encouraged to behave consistently across languages, does its
ranking and reasoning become more reliable?

This question connects directly to the three research questions of the thesis.
RQ1 asks whether optimizing for language-invariant reasoning improves ranking
performance. RQ2 examines how language invariance relates to reasoning quality
across English, Spanish, and Arabic. RQ3 asks whether language-invariant training
objectives lead to more robust multilingual reasoning models. The section first
discusses LLMs as rankers and judges, including the presentation biases that
affect them. It then reviews automated and multilingual fact-checking, before
turning to cross-lingual consistency and the metrics used to measure it. Finally,
it discusses multilingual representations, numerical reasoning across languages,
and the preference-optimization and fine-tuning methods used in this thesis.

\subsection{LLMs as Rankers and Rerankers}
\label{subsec:ranking}

LLMs are now often used as rankers even without task-specific training, especially
in retrieval and evaluation settings. One common approach is pairwise ranking
prompting: the model compares two candidates at a time, and the final ranking is
built by aggregating these pairwise preferences. This makes each decision easier
for the model, but it is expensive because the number of comparisons grows
quadratically with the number of candidates~\cite{qin2024largelanguagemodelseffective}.

Listwise ranking takes a different approach. Instead of comparing candidates in
pairs, the model sees all candidates at once and is asked to produce a complete
ordering. This is more efficient at inference time, but it also makes the output
depend on the full prompt and on the order in which the candidates are shown.
Recent work has adapted learning-to-rank objectives to preference optimization in
order to train listwise rankers more directly~\cite{liu-etal-2025-lipo}. However, these
generative ranking methods remain sensitive to presentation details such as
candidate order and prompt format, rather than depending only on the content of
the candidates~\cite{qiao2026llmbasedlistwisererankingeffect}.

This limitation motivates the independently scored design used in this thesis.
Instead of asking the model to generate a full ranking directly, each candidate
reasoning trace is scored separately and the final ranking is derived from those
scores. This design is especially relevant for the numerical fact-checking task
studied here, where the ranking should reflect the quality of the reasoning trace
rather than artifacts of the prompt.

\subsection{Biases in LLM-as-a-Judge Evaluation}
\label{subsec:judge-bias}

Similar issues appear when LLMs are used as evaluators. Their judgments do not
depend only on the quality of the content being assessed; they can also be shaped
by how the options are presented. Studies of the LLM-as-a-judge setting show
strong position bias, meaning that a candidate's location in the prompt can affect
whether it is preferred. Other effects, such as verbosity bias, have also been
observed~\cite{shi2025judgingjudgessystematicstudy}. These biases matter for ranking because they can
change which candidate is selected, even when the underlying content has not
changed.

For this thesis, the most important issue is that LLM judges are often not
language-invariant. In other words, the same model may give different judgments
when the same evaluation is carried out in different languages. One recent study
examined five model families, five tasks, and twenty-five languages, and found
that multilingual judge agreement was low, with an average Fleiss' $\kappa$ of
roughly $0.3$. The study also found that multilingual pre-training and larger
model scale did not reliably remove this inconsistency~\cite{fu2025reliablemultilingualllmasajudge}.
This kind of cross-lingual instability is exactly the problem that the research
questions in this thesis address.

\subsection{Automated and Multilingual Fact-Checking}
\label{subsec:factcheck}

Automated fact-checking is increasingly treated as a reasoning and assessment
problem for LLMs. Rather than only asking whether a model can produce a label,
recent work studies how the formulation of the task affects the reliability of
model-based verification for real-world claims~\cite{sahitaj2025automatedfactcheckingrealworldclaims}.
Numerical claims make this problem especially difficult. A small change in scope,
units, dates, or quantitative interpretation can be enough to change the correct
verdict.

The QuanTemp benchmark was introduced to study numerical claim verification in a
real-world, open-domain setting, with claims paired with supporting evidence
\cite{v2024quantemprealworldopendomainbenchmark}. Later work builds on this by generating multiple reasoning
traces for each claim and using a verifier to choose the traces that are most
useful for reaching the correct verdict~\cite{chungkham2025thinkrightmoretesttime}. The
CLEF~CheckThat!~2026 Task~2 shared task follows this verifier-guided setup for
multilingual numerical claim verification. In this task, systems work with
English, Spanish, and Arabic claims, rank candidate reasoning traces, and then
derive a claim-level verdict~\cite{clef-checkthat:2026:task2}.

There is also a broader line of work on multilingual fact verification. X-Fact
collects naturally occurring claims across twenty-five languages from eleven
language families, together with expert veracity labels. Results on this benchmark
show that even strong models achieve only modest accuracy, highlighting how
difficult it is to transfer fact-checking ability across languages
\cite{gupta2021xfactnewbenchmarkdataset}. XFEVER extends the FEVER setting by translating claims and
evidence into multiple languages, making it possible to study zero-shot and
translate-train cross-lingual verification directly. It also finds that
performance varies substantially by language and generally remains behind English
\cite{chang2023xfeverexploringfactverification}. Complementing these verification benchmarks, X-FACTR
probes factual knowledge in multilingual models using native-speaker-authored
fill-in-the-blank queries in twenty-three languages, showing that factual recall
also differs across languages~\cite{jiang2020xfactrmultilingualfactualknowledge}.

Together, these studies show why multilingual fact-checking is both important and
difficult. They also provide the practical motivation for this thesis: if
fact-checking systems behave differently across languages, then language
invariance becomes a useful objective to study directly.

\subsection{Language Invariance and Cross-Lingual Consistency}
\label{subsec:invariance}

The main idea behind language invariance is simple: if two inputs express the
same claim in different languages, a robust model should not change its judgment
just because the language changed. This is especially important for the present
task, where the model is expected to reason about the same numerical claim in
English, Spanish, and Arabic. To make this idea useful experimentally, however,
we need to be clear about what kind of consistency is being measured and how it
is separated from accuracy.

\paragraph{Defining consistency.}
A useful starting point is the formal definition of consistency as stability under
meaning-preserving changes to the input. In this view, a model is consistent if it
gives non-contradictory predictions for two quasi-paraphrases of the same query.
The consistency score is then the proportion of paraphrase pairs that lead to
agreeing predictions~\cite{elazar2021measuringimprovingconsistencypretrained}. This definition is important
because it makes consistency separate from accuracy. A model can be consistent and
still be wrong, or it can be accurate on average while giving inconsistent answers
to equivalent inputs.

The same idea naturally extends from paraphrases within one language to
translations across languages. If two versions of a claim have the same meaning,
then a language-invariant model should give the same output for both. Prior work
shows that this is not guaranteed. Multilingual factual-consistency studies find
that encoder models such as mBERT and XLM-R are no more consistent across
languages than a monolingual model is across English paraphrases, and that
inconsistency increases when typologically distant languages are included
\cite{Fierro_2022}. In this thesis, language invariance is therefore treated
as cross-lingual consistency: agreement between a model's outputs for translated
versions of the same input, measured separately from whether those outputs are
correct.

\paragraph{Measuring cross-lingual consistency.}
Because consistency and accuracy are different properties, cross-lingual
consistency should be evaluated on its own. Work on factual knowledge consistency
shows that multilingual models often have low cross-lingual consistency, that
larger model size alone does not reliably solve the problem, and that consistency
is strongly related to subword-vocabulary overlap between languages
\cite{Qi_2023}. This suggests that scripts and tokenization can affect
cross-lingual behavior, not only reasoning ability.

Different output types require different agreement measures. For categorical
outputs such as claim verdicts, standard inter-rater agreement statistics are
appropriate. Cohen's $\kappa$ measures agreement between two raters
\cite{cohen1960coefficient}, Fleiss' $\kappa$ extends this idea to more than two
raters~\cite{fleiss1971measuring}, and Krippendorff's $\alpha$ provides a more
general reliability coefficient that can handle missing data
\cite{krippendorff2004content}. Fleiss' $\kappa$ is also the statistic used in
the multilingual LLM-as-a-judge work discussed earlier~\cite{fu2025reliablemultilingualllmasajudge}.
For ranked outputs, agreement can instead be measured with rank-correlation
statistics such as Kendall's $\tau$ and Spearman's $\rho$
\cite{kendall1938new,spearman1904proof}, or with top-$k$ overlap measures such
as Jaccard similarity.

One caveat is important. Agreement in behavior does not necessarily mean that the
model uses the same internal representation across languages. Recent work shows
that models can answer consistently across languages while still storing
knowledge separately, especially for languages written in different scripts
\cite{ifergan2024beneathsurfaceconsistencyexploring}. For this reason, the thesis treats invariance as an
observable behavioral property, not as proof of a fully shared multilingual
reasoning mechanism.

\paragraph{A measurement framing for this thesis.}
In this thesis, language invariance is measured operationally. For a claim $c$
expressed in the language set
$\mathcal{L} = \{\text{English}, \text{Spanish}, \text{Arabic}\}$, let
$\hat{y}_\ell(c)$ be the verdict predicted from the language-$\ell$ version of
the claim. Our primary metric is \emph{verdict invariance}: whether the predicted
claim verdict remains stable across languages. Concretely, we report the fraction
of aligned claims for which all language-conditioned predictions agree, and we
compute Fleiss' $\kappa$ over the verdicts
$\{\hat{y}_\ell(c)\}_{\ell \in \mathcal{L}}$ by treating the predictions from the
three languages as raters of the same claim~\cite{fleiss1971measuring}. We also
report pairwise Cohen's $\kappa$~\cite{cohen1960coefficient} and Krippendorff's
$\alpha$~\cite{krippendorff2004content} as supporting diagnostics.

The system also produces a ranking of candidate reasoning traces. For a claim
$c$, let $r_\ell(c)$ denote the trace ordering produced from the language-$\ell$
version of the claim. As a secondary diagnostic, we measure ranking-level
invariance by comparing these orderings across languages using mean pairwise
Kendall's $\tau$, Spearman's $\rho$~\cite{kendall1938new,spearman1904proof}, and
top-$k$ Jaccard overlap over the ranked trace indices. These ranking metrics are
not the primary invariance measure, but they help show whether cross-lingual
stability appears in the trace ranking itself or only in the final verdict.

Finally, the training objective uses a different level of invariance. During
training, invariance is encouraged at the score level: the model is penalized when
it assigns different normalized score vectors to the same aligned traces across
languages. This score-level regularizer is related to the consistency objectives
discussed in Section~\ref{subsec:peft}, while the exact formulation is given in
the methodology section. Throughout the thesis, invariance is reported separately
from task accuracy. This separation is essential for RQ2, because a model may be
consistent but wrong, or accurate but unstable across languages
\cite{elazar2021measuringimprovingconsistencypretrained, Qi_2023}.

\subsection{Cross-Lingual Representations and the Latent Language}
\label{subsec:representations}

Language invariance is only plausible if multilingual models share at least some
representation of meaning across languages. A large body of work suggests that
this is partly the case. Studies of mBERT show that it can transfer across
languages in a zero-shot setting, even across different scripts, although transfer
quality depends on factors such as typological similarity~\cite{pires2019multilingualmultilingualbert}.
Other work shows that shared multilingual structure can emerge even without a
shared vocabulary~\cite{wu2020emergingcrosslingualstructurepretrained}. Large-scale multilingual
pre-training, as in XLM-R, further improves transfer across many languages
\cite{conneau2020unsupervisedcrosslingualrepresentationlearning}. Explicit alignment methods can also bring
contextual representations closer across languages and improve cross-lingual
transfer~\cite{cao2020multilingualalignmentcontextualword}. A recent survey provides a broader overview
of these alignment ideas and how they are evaluated~\cite{hämmerl2024understandingcrosslingualalignment}.

For decoder-only LLMs, recent interpretability work points to a related but more
complicated picture. Logit-lens analyses suggest that models such as Llama process
multilingual inputs through intermediate representations that are often closest
to English before producing an answer in the target language. This has been
described as operating through a latent language~\cite{wendler2024llamasworkenglishlatent}.
Similarly, other work finds that semantic decisions in multilingual LLMs are made
in a representation space close to English, and that interventions are most
effective when applied in that space~\cite{schut2025multilingualllmsthinkenglish}. Representational
similarity analyses also suggest that reasoning-like behavior is only partly
aligned across languages~\cite{STOLLE2026100250}.

These findings support one of the assumptions behind this thesis: if multilingual
models already share part of their semantic space, then language-invariant
training may encourage the model to use that shared space more consistently. In
that sense, invariance can act as an inductive bias for more robust multilingual
reasoning (RQ3). However, this should be interpreted carefully. A model can give
the same answer in different languages without necessarily using the same internal
representation in each case~\cite{ifergan2024beneathsurfaceconsistencyexploring}. Therefore, the thesis
treats behavioral invariance as useful evidence of cross-lingual stability, but
not as proof of a fully language-independent reasoning mechanism.

\subsection{Numerical Reasoning Across Languages and Value-Aware Representations}
\label{subsec:numeric}

Numerical reasoning is a known weakness of LLMs. Even when the required operation
is simple, models can behave unreliably. Probing studies describe this as a
fragile number sense, where models fail on elementary numerical tasks that should
in principle be straightforward~\cite{rahman2025fragilenumbersenseprobing}. Smaller diagnostic
tests reach similar conclusions~\cite{rahman2025largelanguagemodelsnumberland}. Benchmarks such as
NumGLUE~\cite{mishra2022numgluesuitefundamentalchallenging} and broader quantitative-reasoning evaluations
\cite{lewkowycz2022solvingquantitativereasoningproblems} show that numerical reasoning is a demanding test of
whether a model understands quantities, rather than only the surrounding text.

One reason this is difficult is that standard token representations do not encode
numerical magnitude directly. A number is usually treated as a sequence of tokens,
and the model has to infer its value from that surface form. Earlier probing work
showed that some embeddings capture numerical magnitude surprisingly well, but
also that subword tokenization limits exact numeracy~\cite{wallace2019nlpmodelsknownumbers}.
Value-aware approaches address this problem by encoding the numerical value of a
span, instead of relying only on its textual representation~\cite{dutulescu2026valueawarenumericalrepresentationstransformer}.
Related number-encoding methods expose canonical numerical information more
directly to the model~\cite{schwartz2024numerologicnumberencodingenhanced}.

This issue becomes even more important in a multilingual setting. The task studied
in this thesis is not only numerical, but also cross-lingual: the model must
handle numerical claims in English, Spanish, and Arabic. Multilingual
chain-of-thought evaluations, including the MGSM benchmark, show that quantitative
reasoning quality varies across languages even when the underlying arithmetic is
the same~\cite{shi2022languagemodelsmultilingualchainofthought}. This means that language can affect numerical
reasoning even when the mathematical problem itself has not changed.

Recent work makes this point even more directly by isolating the surface form of
the number. When the same arithmetic problem is written with different numeral
scripts or formats, model accuracy can drop substantially for underrepresented
scripts. This effect is partly linked to tokenization, although targeted prompting
and explicit numeral mapping can reduce the gap~\cite{reddy2026effectscriptsformatsllm}. Related
analyses show that models also struggle to infer the compositional structure of
numerals across languages unless the relevant operations are made symbolically
explicit~\cite{bhattacharya2025investigatinginteractionlinguisticmathematical}.

These findings matter for this thesis because numeral script is another source of
cross-lingual variation. If a model changes its quantitative judgment when the
same number is written in a different script, then numerical fact-checking becomes
a particularly demanding setting for language invariance. Value-aware numerical
representations and numerical normalization are therefore relevant not only as a way to improve numerical
accuracy, but also as a possible way to reduce cross-lingual inconsistency. This
is especially important for the Arabic part of the task, where script and numeral
format may introduce additional variation.

\subsection{Preference Optimization and Parameter-Efficient Fine-Tuning}
\label{subsec:peft}

The training setup in this thesis builds on two practical lines of work:
preference optimization and parameter-efficient fine-tuning. Direct Preference
Optimization (DPO) trains a model from pairwise preferences by increasing the
relative likelihood of preferred outputs over rejected ones, while keeping the
model close to a reference policy. This makes it possible to optimize preferences
without training a separate reward model or using reinforcement learning
\cite{rafailov2024directpreferenceoptimizationlanguage}.

For the fine-tuning itself, parameter-efficient methods are important because the
models used in this thesis are too large to update fully under a modest GPU
budget. Low-Rank Adaptation (LoRA) addresses this by freezing the base model and
training only a small set of low-rank update matrices~\cite{hu2021loralowrankadaptationlarge}.
Quantized LoRA (QLoRA) reduces memory use further by loading the base model in
low precision and training adapters on top of it~\cite{dettmers2023qloraefficientfinetuningquantized}. These
methods make it feasible to adapt open-weight instruction-tuned models such as
Mistral-7B~\cite{jiang2023mistral7b}, Qwen3-8b\cite{qwen3technicalreport} and Llama-3.1-8B-Instruct
\cite{llama31instruct}, which are the scale of models used in the experiments.

Recent work has also started to use these kinds of objectives for cross-lingual
consistency. Consistency regularization penalizes a model when its predictions
change under meaning-preserving perturbations, such as code-switching or machine
translation~\cite{zheng2021consistencyregularizationcrosslingualfinetuning}. Cross-lingual self-alignment extends the
preference-optimization idea by sampling answers in several languages, selecting
the most self-consistent answer as the preferred target, and applying DPO to align
the model's knowledge across languages~\cite{wang2025calmunleashingcrosslingualselfaligning}. Direct Consistency
Optimization is even closer to the goal of this thesis: it adapts the DPO
formulation to align completion likelihoods across parallel prompts in different
languages, without requiring an explicit reward model~\cite{liu2026posttraininglanguagemodelscrosslingual}.

These methods show that preference optimization can be used for consistency, not
only for general output quality. However, most of this work evaluates consistency
through factual question answering or broad benchmarks. It does not directly test
whether optimizing for cross-lingual consistency improves ranking performance in
a numerical fact-checking task. That is the gap addressed by RQ1: whether an
invariance-oriented objective can improve the ranking itself, rather than simply
making outputs agree across languages.

\subsection{Positioning of This Thesis}
\label{subsec:positioning}

The work reviewed above points to a clear gap. LLM rankers and judges are
sensitive to presentation details such as candidate order and prompt format
\cite{shi2025judgingjudgessystematicstudy, qiao2026llmbasedlistwisererankingeffect}.
They are also often not language-invariant: the same model can produce different
verdicts or rankings for semantically equivalent inputs in different languages,
and this inconsistency is not reliably solved by larger model scale or
multilingual pre-training~\cite{fu2025reliablemultilingualllmasajudge, Qi_2023,
Fierro_2022}. At the same time, multilingual models do appear to learn partly
shared, partly English-centric representations~\cite{wendler2024llamasworkenglishlatent,
conneau2020unsupervisedcrosslingualrepresentationlearning}. This suggests that
language-invariant behavior may be attainable, while still leaving open the
caution that behavioral agreement does not necessarily prove a shared internal
representation~\cite{ifergan2024beneathsurfaceconsistencyexploring}.

The numerical side of the task adds another layer of difficulty. Prior work shows
that numerical reasoning is fragile and only partly captured by standard token
representations~\cite{wallace2019nlpmodelsknownumbers,
rahman2025fragilenumbersenseprobing}. It is also sensitive to numeral script and
format across languages~\cite{reddy2026effectscriptsformatsllm}. This matters for
the present thesis because the task is not only multilingual, but also numerical:
the model must rank reasoning traces for claims whose quantities may appear in
different linguistic and script forms. Finally, recent work shows that preference
optimization and consistency regularization can be aimed at cross-lingual
consistency rather than output quality alone
\cite{liu2026posttraininglanguagemodelscrosslingual,
wang2025calmunleashingcrosslingualselfaligning,
zheng2021consistencyregularizationcrosslingualfinetuning}. What remains unclear
is whether treating language invariance as a training objective actually improves
multilingual reasoning in a trace-ranking fact-checking setting.

This thesis addresses that gap using the CLEF~CheckThat!~2026 Task~2 setting,
where systems rank candidate reasoning traces and produce a final verdict for
numerical claims in English, Spanish, and Arabic. The three research questions
separate the problem into ranking performance, the relationship between invariance
and reasoning quality, and broader multilingual robustness.

\paragraph{RQ1: Does optimizing for language-invariant reasoning improve ranking performance?}
Prior work shows that consistency can be optimized directly
\cite{zheng2021consistencyregularizationcrosslingualfinetuning,
wang2025calmunleashingcrosslingualselfaligning,
liu2026posttraininglanguagemodelscrosslingual}, but most of this work evaluates
the result on factual question answering or general benchmarks. RQ1 asks whether
the same idea helps in a ranking-based fact-checking task. In particular, it
tests whether adding a language-invariance objective to the trace scorer improves
the official ranking metrics, rather than only making outputs more similar across
languages. This connects to work on ranking bias
\cite{shi2025judgingjudgessystematicstudy,
qiao2026llmbasedlistwisererankingeffect,
qin2024largelanguagemodelseffective}: if rankings are affected by irrelevant
variation in the input, then reducing language-driven variation may make the
ranking more reliable.

\paragraph{RQ2: What is the relationship between language invariance and reasoning quality across languages?}
Consistency is not the same as accuracy~\cite{elazar2021measuringimprovingconsistencypretrained,
Qi_2023}. A model can be stable across languages but wrong, or accurate in one
language while unstable across others. RQ2 therefore measures invariance and
per-language reasoning quality separately, then analyzes how they relate across
English, Spanish, and Arabic. The thesis does not assume that the relationship is
simple or monotonic. Prior work suggests that cross-lingual consistency can be
affected by vocabulary overlap and script~\cite{Qi_2023}, and that agreement
across languages does not necessarily imply shared internal representations
\cite{ifergan2024beneathsurfaceconsistencyexploring}. For this reason, RQ2 treats
the link between invariance and reasoning quality as an empirical question.

\paragraph{RQ3: Can language-invariant training objectives lead to more robust multilingual reasoning models?}
The representational evidence reviewed above suggests that multilingual models
share part of their semantic space
\cite{wendler2024llamasworkenglishlatent, schut2025multilingualllmsthinkenglish,
conneau2020unsupervisedcrosslingualrepresentationlearning}. RQ3 asks whether a
language-invariant objective can use this shared structure as an inductive bias
for more robust multilingual reasoning. This is tested across English, Spanish,
and Arabic, with particular attention to Arabic because script and tokenization
can introduce additional variation~\cite{reddy2026effectscriptsformatsllm}. The
numerical dimension is also important: value-aware numerical representations and
number-encoding methods are considered as complementary ways to reduce variation
caused by the surface form of quantities
\cite{dutulescu2026valueawarenumericalrepresentationstransformer,
schwartz2024numerologicnumberencodingenhanced}.